%% ma: L12-B15-0029
%% lop: 12
%% chuong: 5
%% bai: 15
%% dang: 12.15.11
%% loai: TN4
%% muc: NB
%% dap_an: "B"
%% nguon: "To_n_12_PT_M_t_Ph_ng.pdf (Đề kiểm tra đánh giá kiến thức, Chuyên đề Phương trình mặt phẳng, Nguyễn Anh Khôi - Lê Minh Khôi), Câu 10"
%% kien_thuc: "Bài 15 (vectơ chỉ phương từ phương trình chính tắc, đường thẳng vuông góc với mặt phẳng), Bài 14 (vectơ pháp tuyến)"
%% trang_thai: chinh-thuc
%% ghi_chu: "Tài liệu gốc không kèm đáp án; đáp án do Claude giải và kiểm bằng sympy. Trong tài liệu chuyên đề Phương trình mặt phẳng nhưng dùng phương trình đường thẳng nên xếp vào Bài 15."
\begin{ex}
Trong không gian $Oxyz$, đường thẳng $\Delta\colon \dfrac{x}{1}=\dfrac{y}{1}=\dfrac{z}{2}$ vuông góc với mặt phẳng nào trong các mặt phẳng sau?
\choice{$(P)\colon x+y+z=0$}{\True $(\alpha)\colon x+y+2z=0$}{$(\beta)\colon x+y-z=0$}{$(Q)\colon x+y-2z=0$}
\loigiai{$\Delta$ có vectơ chỉ phương $\vec u=(1;1;2)$. $\Delta$ vuông góc với mặt phẳng nhận $\vec u$ làm vectơ pháp tuyến, đó là $(\alpha)\colon x+y+2z=0$. Chọn B.}
\end{ex}
